\Needspace{18\baselineskip}
\section{Exact inventories and sharp context closure at every depth}
\label{sec:deep-inventory}

The provenance argument establishes a useful upper bound on hole inventories.
We now strengthen it to an equality described entirely by the original host and
the selected address. This also gives a necessary-and-sufficient freshness
condition. Throughout this section $s$ is a site in $G$, $H$ has the interface
and local depth required by $s$, and $G[s:=H]$ denotes $\mathsf{plugAt}(s,H)$.
Write $\mathcal H(G)$ for the set of full typed hole labels present in $G$.
These statements concern the concrete typed hierarchy; they do not presuppose
any specimen laws or a concrete UAM semantics.

\subsection{Separating one outer edge from the rest of a graph}

\begin{definition}[Inventory of an edge]\label{term:edge-hole-inventory}
For an outer edge $e$ of $G$, let $E_G(e,h)$ mean that some position in the fibre
of $e$ observes $h$. This includes the edge itself and all descendants within
its label. Thus $E_G(e,h)$ is the predicate $\mathsf{EdgeHasHole}(G,e,h)$.
\end{definition}
\lean{Graph/DeepInventory}{9}{10}{Linear.EdgeHasHole}

\begin{definition}[Uniqueness away from an edge]\label{term:away-unique}
The predicate $\mathsf{AwayUnique}(G,e)$ says that any two occurrences with outer
edges different from $e$ and with the same hole observation are equal. It places
no condition on occurrences within the fibre of $e$.
\end{definition}
\lean{Graph/DeepInventory}{13}{15}{Linear.AwayUnique}

\begin{lemma}[Partition of uniqueness]
The graph $G$ has unique holes if and only if all three conditions hold:
occurrences away from $e$ are unique; hole observations within the fibre of $e$
are unique; and no label is observed both away from $e$ and within that fibre.
\end{lemma}
\begin{proof}
Global uniqueness restricts to each part. A collision between the two parts
would force equality of their outer edges, contradicting the requirement that
the first outer edge differ from $e$. Conversely, take two occurrences of the
same hole label. If both lie away from $e$, use away-uniqueness. If both have
outer edge $e$, use fibre uniqueness. The two mixed cases contradict the
cross-part disjointness condition. These cases exhaust all occurrences.
\end{proof}
\lean{Graph/DeepInventory}{19}{43}{Linear.uniqueHoles_partition}

\begin{lemma}[Inventories of a box and of a whole graph]
If the label of $e$ is a box with body $B$, then $E_G(e,h)$ holds exactly when
$h\in\mathcal H(B)$. For an arbitrary edge $e$, a label $h$ occurs in $G$ exactly
when it occurs away from $e$ or satisfies $E_G(e,h)$.
\end{lemma}
\begin{proof}
The shallow position of a box is not a hole, while each of its child positions
is precisely an occurrence of its body, with the same observation. For the
whole-graph statement, write an occurrence as an outer edge and a fibre
position, and split according to whether that edge equals $e$. In the reverse
direction either witness is already an occurrence of $G$.
\end{proof}
\lean{Graph/DeepInventory}{46}{56}{Linear.edgeHasHole_box}
\lean{Graph/DeepInventory}{59}{69}{Linear.hasHole_partition}

\begin{lemma}[Exact inventory after replacing a box body]
Changing the body of box edge $e$ to $B'$ preserves the inventory away from $e$
exactly. The entire new inventory is the union of the old inventory away from
$e$ and $\mathcal H(B')$.
\end{lemma}
\begin{proof}
Every other edge has exactly the same label before and after the change.
Transport its fibre positions along this label equality; fibre transport
preserves observations in both directions. This proves the first claim.
Partition the new graph at $e$ and use the preceding box-inventory lemma to
identify its selected fibre with the new body's inventory.
\end{proof}
\lean{Graph/DeepInventory}{72}{84}{Linear.replaceBox_retained_iff}
\lean{Graph/DeepInventory}{87}{92}{Linear.replaceBox_hasHole_iff}

\begin{lemma}[Box uniqueness]
The hole observations in a box fibre are unique exactly when its body has
unique holes. Consequently, the body of every box in a context has unique holes.
\end{lemma}
\begin{proof}
A box's shallow position cannot witness a hole. Two child positions observing
the same hole are equal precisely when their body occurrences are equal.
Injectivity of the child constructor proves the reverse implication. For the
consequence, restrict the graph's uniqueness to the selected fibre using the
partition lemma, then apply this equivalence.
\end{proof}
\lean{Graph/DeepInventory}{95}{109}{Linear.boxFiber_unique_iff}
\lean{Graph/DeepInventory}{112}{117}{Linear.uniqueHoles_box_body}

\begin{lemma}[Uniqueness after changing one box]
Suppose $G$ is a context. Replacing the body of box edge $e$ by $B'$ yields a
context if and only if $B'$ has unique holes and its hole inventory is disjoint
from the inventory of $G$ away from $e$.
\end{lemma}
\begin{proof}
First the unchanged part remains unique: transport two of its occurrences to
$G$ along their unchanged label equalities, apply uniqueness in $G$, and use
injectivity of fibre transport to recover equality in the new graph.
The uniqueness partition for the changed graph now has a proved first
condition. Its second condition is equivalent to uniqueness of $B'$, by box
uniqueness. Its third condition is exactly the stated disjointness, by the
preservation of the outside inventory and the exact inventory of a box.
\end{proof}
\lean{Graph/DeepInventory}{120}{133}{Linear.replaceBox_awayUnique}
\lean{Graph/DeepInventory}{137}{157}{Linear.replaceBox_uniqueHoles_iff}

\subsection{A syntactic inventory of retained holes}

\begin{definition}[Retention along a site]\label{term:site-retained-hole}
Let $R_s(h)$ be defined recursively on the site address. At a shallow site
selecting $e$ in $G$, it is $\mathsf{RetainedHole}(G,e,h)$: there is an occurrence
of $h$ under another outer edge. If the site descends through box edge $e$ to an
inner site $t$, set
\[
 R_s(h)\quad\Longleftrightarrow\quad
 \mathsf{RetainedHole}(G,e,h)\ \lor\ R_t(h).
\]
This is $\mathsf{SiteRetainedHole}$. It depends only on the host and address,
not on the replacement graph or a provenance certificate. It excludes all
occurrences below the edge finally removed, while retaining the siblings of
every enclosing box. Enclosing box edges are not themselves holes.
\end{definition}
\lean{Graph/DeepInventory}{161}{165}{Linear.SiteRetainedHole}

\begin{theorem}[Exact deep hole inventory]
For every site and every type-correct replacement, without uniqueness or
freshness assumptions,
\[
 h\in\mathcal H(G[s:=H])\quad\Longleftrightarrow\quad
 R_s(h)\ \lor\ h\in\mathcal H(H).
\]
\end{theorem}
\lean{Graph/DeepInventory}{168}{176}{Linear.plugAt_hasHole_iff}
\begin{proof}
Induct on the site address. The shallow case is the exact shallow inventory
theorem. For a site inside a box, the exact box-replacement inventory is the
union of the unaffected outer inventory and the changed body's inventory.
The induction hypothesis identifies the latter as the inner retained inventory
union the replacement inventory. Reassociate these three alternatives; the
first two are exactly $R_s$. Both directions are preserved at each step, so no
retained or inserted label is lost from the characterization.
\end{proof}

\begin{definition}[Exact deep freshness]\label{term:site-fresh}
The predicate $\mathsf{SiteFresh}(s,H)$ says
\[
 \forall h,\quad R_s(h)\ \Longrightarrow\ h\notin\mathcal H(H).
\]
Only retained host labels constrain the replacement; a label that is wholly
consumed is available for reuse.
\end{definition}
\lean{Graph/DeepInventory}{179}{181}{Linear.SiteFresh}

\begin{lemma}[Retention implies original presence]
If $R_s(h)$, then $h\in\mathcal H(G)$.
\end{lemma}
\begin{proof}
Induct on the site. An outside witness is already a host occurrence. For an
inner retained witness, the induction hypothesis supplies an occurrence in the
original box body; the box-inventory lemma and whole-graph partition lift it
back into the host.
\end{proof}
\lean{Graph/DeepInventory}{184}{193}{Linear.siteRetainedHole_hasHole}

\begin{lemma}[Replacement uniqueness is necessary for shallow plugging]
If a shallow plugged graph has unique holes, its replacement has unique holes.
\end{lemma}
\begin{proof}
The right-tagged injection embeds every replacement edge into the result,
preserving labels and all fibre observations. Uniqueness reflects along this
injective edge map by the established embedding theorem.
\end{proof}
\lean{Graph/DeepInventory}{196}{200}{Linear.plug_uniqueHoles_replacement}

\subsection{Necessary and sufficient context closure}

\begin{theorem}[Exact context criterion at arbitrary depth]
Suppose $G$ has unique holes. Then
\[
 \mathsf{UniqueHoles}(G[s:=H])
 \quad\Longleftrightarrow\quad
 \mathsf{UniqueHoles}(H)\ \land\ \mathsf{SiteFresh}(s,H).
\]
Thus uniqueness of the replacement is a necessary conclusion in the forward
direction, not an additional assumption of the theorem.
\end{theorem}
\begin{proof}
Induct on the address. At a shallow site, uniqueness of the result forces both
replacement uniqueness and sharp shallow freshness. The reverse implication is
the proved shallow closure theorem.

For an inside site, let $B$ be the old body and $t$ the inner address. Since $G$
is a context, $B$ is a context. The box-replacement criterion says that the
result is a context exactly when $B[t:=H]$ is a context and its inventory is
disjoint from the outer inventory away from the selected box. The induction
hypothesis replaces its first condition by uniqueness of $H$ and inner
freshness.

For necessity, inner freshness excludes collisions with inner retained labels.
A label inserted from $H$ is present in the changed body by the exact inventory
theorem; outer disjointness therefore excludes collisions with every retained
outer label. These are precisely the two alternatives defining $R_s$.

For sufficiency, site freshness supplies inner freshness, so induction proves
uniqueness of the changed body. Consider a label shared by that body and an
unaffected outer edge. The exact inventory theorem gives two cases. If the
label comes from $H$, site freshness is contradicted. If it is retained inside
the body, retention implies its presence in the old body. It would then have
occurred both inside and outside that box already in $G$, contradicting the
uniqueness partition of $G$. This proves the required outer disjointness and
completes the induction.
\end{proof}
\lean{Graph/DeepInventory}{204}{234}{Linear.plugAt_uniqueHoles_iff}

If $H$ is already a context, its uniqueness conjunct is known, leaving the
necessary-and-sufficient condition $\mathsf{SiteFresh}(s,H)$ alone.
\lean{Graph/DeepInventory}{237}{242}{Linear.plugAt_uniqueHoles_iff_fresh}
\termdef{deep-certified-fresh}
The constructor $\mathsf{Context.plugAtFresh}$ pairs the deep replacement graph
with the proof supplied by this equivalence. Unlike full-label disjointness,
its required certificate permits reuse of consumed labels at any box depth.
\lean{Graph/DeepInventory}{245}{248}{Linear.Context.plugAtFresh}

\begin{theorem}[Exact hole-freeness criterion]
The result $G[s:=H]$ is hole-free if and only if no label satisfies $R_s$ and $H$
is hole-free. No uniqueness assumption is needed.
\end{theorem}
\begin{proof}
The exact inventory is empty precisely when both of its constituent inventories
are empty. In either direction, an alleged witness in one part gives, or comes
from, a witness in the union.
\end{proof}
\lean{Graph/DeepInventory}{252}{260}{Linear.plugAt_holeFree_iff}

\subsection{Consumption of the selected label}

\begin{definition}[Selected hole label]\label{term:selected-hole}
The partial observation $\mathsf{Site.selectedHole}(s)$ follows the address to
its selected edge and reads that edge's own shallow hole label. It does not
inspect descendants of the selected edge. Thus selecting a box yields no hole
label, even if its body contains holes.
\end{definition}
\lean{Graph/DeepInventory}{263}{267}{Linear.Site.selectedHole}

\begin{lemma}[Presence and non-retention of a selected hole]
If the selected label is $h$, then $h$ occurs in the original host. If the host
is a context, $R_s(h)$ is false.
\end{lemma}
\begin{proof}
For presence, the shallow case uses the selected edge's shallow fibre position.
The inside case lifts the induction hypothesis from the body through the box.
For non-retention, the shallow statement is the existing consumed-label lemma.
Inside a box, a retained witness outside the box contradicts the old graph's
uniqueness partition, because the selected label occurs inside its body. An
inner retained witness contradicts induction applied to that body's uniqueness.
\end{proof}
\lean{Graph/DeepInventory}{270}{276}{Linear.Site.selectedHole_present}
\lean{Graph/DeepInventory}{280}{291}{Linear.site_selected_not_retained}

\begin{theorem}[Exact consumption at arbitrary depth]
If $G$ is a context and $s$ selects a hole labelled $h$, then
\[
 h\in\mathcal H(G[s:=H])\quad\Longleftrightarrow\quad h\in\mathcal H(H).
\]
\end{theorem}
\begin{proof}
Use the exact deep inventory theorem and eliminate its retained-host alternative
by non-retention of the selected label. This permits both genuine removal and
intentional reintroduction, depending on the replacement.
\end{proof}
\lean{Graph/DeepInventory}{294}{299}{Linear.plugAt_consumed_hole_iff}

\subsection{Concrete regressions inside boxes}

In the existing one-hole boxed graph, no label is retained when the inner hole
is selected. At the outer level there is no other edge (the carrier has one
element), and inside the box there is no other edge either. This proves
$\neg R_s(h)$ for every $h$.
\lean{Graph/DeepInventoryExamples}{9}{14}{Linear.Examples.boxedHole_retained_absent}
Replacing that hole by the same one-hole context therefore preserves uniqueness,
by sharp deep freshness. Full disjointness fails because both original graphs
contain the selected label. The consumed-label theorem also proves that this
label is genuinely present after replacement, witnessed by its occurrence in
the inserted graph.
\lean{Graph/DeepInventoryExamples}{17}{19}{Linear.Examples.deep_reuse_unique}
\lean{Graph/DeepInventoryExamples}{21}{23}{Linear.Examples.deep_reuse_not_disjoint}
\lean{Graph/DeepInventoryExamples}{26}{27}{Linear.Examples.deep_reuse_present}
Replacing instead with the hole-free value graph yields a hole-free result:
combine absence of retained labels and hole-freeness of the value using the
exact hole-freeness criterion. This is a shorter independent derivation of the
previous global filling result.
\lean{Graph/DeepInventoryExamples}{30}{31}{Linear.Examples.deep_fill_holeFree}

For a negative example, box the earlier two-hole chain with zero bound variables
and one free output. Its external interface is one zero-binding thunk input and
one term output. Define a site descending into this box and selecting the first
hole. The boxed graph is a context because the two-hole body is a context and
boxing preserves uniqueness.
\lean{Graph/DeepInventoryExamples}{34}{35}{Linear.Examples.boxedTwoHoles}
\lean{Graph/DeepInventoryExamples}{37}{38}{Linear.Examples.boxedFirstSite}
\lean{Graph/DeepInventoryExamples}{40}{41}{Linear.Examples.boxedTwoHoles_unique}
The one-hole replacement carrying label~$1$ is also a context, but it collides
with the retained second hole. These two witnesses refute deep freshness. The
necessary direction of the deep context theorem then proves that the result
fails uniqueness.
\lean{Graph/DeepInventoryExamples}{44}{46}{Linear.Examples.deep_collision_not_fresh}
\lean{Graph/DeepInventoryExamples}{48}{53}{Linear.Examples.deep_collision_breaks_unique}
Finally, the exact inventory distinguishes the labels: label~$0$ is absent and
label~$1$ is present. The first claim uses consumption and the fact that the
replacement's label has name~$1$, not~$0$; the second uses its inserted witness.
An inventory records presence rather than multiplicity, so this is consistent
with the independently proved duplicate occurrence of label~$1$.
\lean{Graph/DeepInventoryExamples}{56}{65}{Linear.Examples.deep_collision_inventory}
